\documentclass[11pt,a4paper]{article}
\usepackage[T1]{fontenc}
\usepackage[utf8]{inputenc}
\usepackage[polish]{babel}
\usepackage{lmodern}
\usepackage[a4paper,margin=24mm,headheight=15pt]{geometry}
\usepackage{amsmath,amssymb}
\usepackage{graphicx}
\usepackage{booktabs,tabularx}
\usepackage{xcolor}
\usepackage{listings}
\usepackage{microtype}
\usepackage{fancyhdr}
\usepackage{caption}
\usepackage[section]{placeins}
\usepackage{hyperref}

\definecolor{accent}{HTML}{1D4ED8}
\definecolor{codebg}{HTML}{F3F6FA}
\hypersetup{colorlinks=true,linkcolor=accent,urlcolor=accent,citecolor=accent}
\captionsetup{font=small,labelfont=bf}
\setlength{\parindent}{0pt}
\setlength{\parskip}{6pt}
\setlength{\emergencystretch}{2em}
\pagestyle{fancy}
\fancyhf{}
\fancyhead[L]{\small Algorytmy macierzowe}
\fancyhead[R]{\small Laboratorium 1}
\fancyfoot[C]{\thepage}
\lstset{basicstyle=\ttfamily\footnotesize,backgroundcolor=\color{codebg},
  keywordstyle=\color{accent}\bfseries,frame=single,rulecolor=\color{black!12},
  breaklines=true,columns=fullflexible,keepspaces=true,showstringspaces=false,
  xleftmargin=4pt,xrightmargin=4pt,aboveskip=9pt,belowskip=9pt}

\newcommand{\autorzy}{Imię i nazwisko 1, Imię i nazwisko 2}
\newcommand{\grupa}{Numer grupy}
\newcommand{\dataoddania}{Data oddania}
\newcommand{\wykres}[3]{
  \begin{figure}[htbp]
    \centering
    \IfFileExists{plots/#1}{\includegraphics[width=0.97\linewidth]{plots/#1}}{
      \fbox{\parbox[c][45mm][c]{0.9\linewidth}{\centering
      Brak pliku wykresu: \texttt{\detokenize{#1}}}}}
    \caption{#2}\label{#3}
  \end{figure}}

\begin{document}

\begin{center}
  {\small Akademia Górniczo-Hutnicza w Krakowie}\par
  \vspace{8mm}
  {\Large\bfseries Algorytmy macierzowe}\par
  \vspace{3mm}
  {\LARGE\bfseries Rekurencyjne mnożenie macierzy}\par
  \vspace{3mm}
  {\large Binét, Strassen i AlphaTensor}\par
  \vspace{7mm}
  \autorzy\par
  Grupa: \grupa\quad |\quad \dataoddania
\end{center}
\vspace{5mm}

\begin{tabularx}{\linewidth}{@{}lX@{}}
\toprule
Procesor & AMD Ryzen 5 5500, 6 rdzeni, 12 wątków \\
Pamięć RAM & 16 GB \\
Karta graficzna & AMD Radeon RX 6700 XT, 12 GB VRAM \\
System & Windows 11 Home, 64-bit \\
Kompilacja & Python 3.14.4, Cython 3.3.0, MSVC 14.50; moduły \texttt{.pyd} \\
Weryfikacja & GNU Octave; dodatkowo NumPy \\
\bottomrule
\end{tabularx}

\vspace{4mm}
\textbf{Status danych.} Implementacje zostały skompilowane i przetestowane.
Pełne pomiary modułów Cython nie zostały jeszcze wykonane.
Wykresy i tabele w tej wersji przedstawiają wcześniejsze pomiary zwykłych
modułów Python. Parametry środowiska powyżej opisują komputer i aktualną
konfigurację kompilacji; nie stanowią potwierdzenia wersji środowiska
każdego wcześniejszego pomiaru.

\tableofcontents
\newpage

\section{Cel i zakres ćwiczenia}
Celem ćwiczenia jest implementacja rekurencyjnych metod mnożenia macierzy,
weryfikacja poprawności oraz porównanie czasu, liczby operacji arytmetycznych
i szczytu śledzonych alokacji pamięci.
Wspólne testy wydajności dotyczą macierzy kwadratowych $n\times n$.
Binét i AlphaTensor dodatkowo obsługują zgodne macierze prostokątne.

Zaimplementowano algorytmy Binéta, Strassena oraz metodę wykorzystującą
faktoryzację AlphaTensor dla iloczynu $4\times5$ przez $5\times5$.
Pierwotna treść zestawu wymienia w punkcie trzecim metodę Ladermana.
AlphaTensor stanowi wybrany wariant tego punktu; zakres tej zamiany
należy uzgodnić z prowadzącym przed oddaniem sprawozdania.

Elementy macierzy wejściowych losowane są z przedziału otwartego
$(10^{-8},1)$ i przechowywane jako liczby zmiennoprzecinkowe
podwójnej precyzji. Skrajne wartości przedziału wykluczono przez
\texttt{np.nextafter}. Macierze do wspólnych porównań zapisano w CSV,
stosując generator PCG64 i ziarno 42. Kolejne metody wczytują te same
pliki zamiast ponownie losować dane.

\section{Implementacja i kompilacja}
Kod algorytmów zapisano w składni Python. Cython tłumaczy pliki źródłowe
na C, a kompilator MSVC tworzy natywne moduły rozszerzeń \texttt{.pyd}.
Skrypty uruchamiające pomiary i tworzące wykresy pozostają skryptami Python.
Nie twierdzimy zatem, że Python jako język przestał być interpretowany;
skompilowane są moduły realizujące badane algorytmy.

\begin{lstlisting}[language=Python,caption={Przykład konfiguracji kompilacji Binéta.}]
setup(
    name="binet_compiled",
    ext_modules=cythonize(
        [Extension("binet_compiled", ["binet.py"])],
        compiler_directives={"language_level": 3},
        build_dir="build/cython",
    ),
)
\end{lstlisting}

\begin{lstlisting}[caption={Polecenia budowania trzech modułów.}]
python setup_binet.py build_ext --inplace
python setup_strassen.py build_ext --inplace
python setup_alphatensor.py build_ext --inplace
\end{lstlisting}

Moduł AlphaTensor korzysta ze skompilowanej funkcji
\texttt{multiply\_direct} z modułu \texttt{strassen\_compiled}.
Kompilacja nie zmienia wzorów ani liczby zliczanych operacji.
Nie zastosowano dodatkowych deklaracji typów C dla wszystkich zmiennych,
więc część operacji nadal korzysta z obiektów Python i biblioteki NumPy.
Sam fakt kompilacji nie gwarantuje określonego przyspieszenia.
Obliczenia nie wykorzystują GPU.

\section{Model zliczania operacji}
Każde mnożenie dwóch skalarów oraz każde binarne dodawanie lub odejmowanie
liczymy jako jedną operację. Funkcje zwracają trójkę
$(C,M,S)$: wynik, liczbę mnożeń oraz liczbę dodawań i odejmowań.
Łączna liczba operacji wynosi
\[
  F=M+S.
\]
Nie wykonujemy dzieleń zmiennoprzecinkowych w badanych algorytmach.
Nie zliczamy przypisań, indeksowania, alokacji, kopiowania ani operacji
całkowitoliczbowych na indeksach. W AlphaTensor nie liczymy osobno
jednoargumentowej zmiany znaku; binarne odejmowanie jest zliczane.
Operacje na zerach dodanych przy dopełnianiu macierzy są uwzględniane.

Dodanie dwóch bloków o rozmiarze $r\times s$ oznacza $rs$ dodawań
liczb. Dlatego po złożeniu wyniku Binéta sumujemy także rozmiary
bloków $C_{11},C_{12},C_{21},C_{22}$.

\begin{lstlisting}[language=Python,caption={Zliczanie dodawań w jednym kroku Binéta.}]
additions = (
    add1 + add2 + add3 + add4
    + add5 + add6 + add7 + add8
    + C11.size + C12.size + C21.size + C22.size
)
\end{lstlisting}

\section{Algorytm Binéta}
\subsection{Podział na bloki i pseudokod}
Dla zgodnych podziałów macierzy
\[
A=\begin{bmatrix}A_{11}&A_{12}\\A_{21}&A_{22}\end{bmatrix},\qquad
B=\begin{bmatrix}B_{11}&B_{12}\\B_{21}&B_{22}\end{bmatrix}
\]
wyznaczamy osiem iloczynów rekurencyjnie i tworzymy cztery bloki wyniku:
\begin{align*}
C_{11}&=A_{11}B_{11}+A_{12}B_{21}, &
C_{12}&=A_{11}B_{12}+A_{12}B_{22},\\
C_{21}&=A_{21}B_{11}+A_{22}B_{21}, &
C_{22}&=A_{21}B_{12}+A_{22}B_{22}.
\end{align*}

\begin{lstlisting}[caption={Pseudokod Binéta dla zgodnych macierzy prostokątnych.}]
Binet(A, B):
    (r, k) = rozmiar(A)
    s = liczba_kolumn(B)
    if r == 1 or k == 1 or s == 1:
        return mnozenie_bezposrednie(A, B)
    podziel A i B w polowie kazdego wymiaru
    C11 = Binet(A11, B11) + Binet(A12, B21)
    C12 = Binet(A11, B12) + Binet(A12, B22)
    C21 = Binet(A21, B11) + Binet(A22, B21)
    C22 = Binet(A21, B12) + Binet(A22, B22)
    zsumuj liczniki wywolan oraz dodawan blokow
    return polacz(C11, C12, C21, C22), M, S
\end{lstlisting}

\subsection{Dowolne rozmiary}
Każdy wymiar dzielimy na części o długościach $\lfloor d/2\rfloor$
i $\lceil d/2\rceil$. Nie dopełniamy macierzy zerami.
Wymiar wspólny jest dzielony w tym samym miejscu dla kolumn $A$
i wierszy $B$, co zachowuje zgodność iloczynów blokowych.
Gdy którykolwiek wymiar osiągnie 1, wykonujemy mnożenie bezpośrednie.

\begin{lstlisting}[language=Python,caption={Podział o nierównych rozmiarach bloków.}]
mid_rows = rows // 2
mid_common = common // 2
mid_cols = cols // 2
A11 = A[:mid_rows, :mid_common]
A12 = A[:mid_rows, mid_common:]
B11 = B[:mid_common, :mid_cols]
B21 = B[mid_common:, :mid_cols]
P1, mul1, add1 = binet(A11, B11)
P2, mul2, add2 = binet(A12, B21)
C11 = P1 + P2
\end{lstlisting}

\subsection{Złożoność}
Dla $A\in\mathbb{R}^{r\times k}$, $B\in\mathbb{R}^{k\times s}$
zastosowana implementacja wykonuje dokładnie
\[
 M=rks,\qquad S=rs(k-1).
\]
Dla kwadratowych macierzy $M=n^3$, $S=n^2(n-1)$,
a $F=2n^3-n^2$. Rekurencja dla $n=2^j$ ma postać
$T(n)=8T(n/2)+\Theta(n^2)$, więc $T(n)=\Theta(n^3)$.

\section{Algorytm Strassena}
\subsection{Siedem iloczynów zamiast ośmiu}
Dla parzystego $n$ dzielimy obie macierze na równe bloki i wyznaczamy
\begin{align*}
P_1&=(A_{11}+A_{22})(B_{11}+B_{22}),\\
P_2&=(A_{21}+A_{22})B_{11}, & P_3&=A_{11}(B_{12}-B_{22}),\\
P_4&=A_{22}(B_{21}-B_{11}), & P_5&=(A_{11}+A_{12})B_{22},\\
P_6&=(A_{21}-A_{11})(B_{11}+B_{12}),\\
P_7&=(A_{12}-A_{22})(B_{21}+B_{22}).
\end{align*}
Bloki wyniku obliczamy następująco:
\begin{align*}
C_{11}&=P_1+P_4-P_5+P_7, & C_{12}&=P_3+P_5,\\
C_{21}&=P_2+P_4, & C_{22}&=P_1-P_2+P_3+P_6.
\end{align*}

\begin{lstlisting}[caption={Pseudokod Strassena dla dowolnego dodatniego rozmiaru kwadratowego.}]
Strassen(A, B):
    n = liczba_wierszy(A)
    if n == 1:
        return A * B, 1, 0
    if n jest nieparzyste:
        return StrassenOdd(A, B)
    podziel A i B na cztery rowne bloki
    oblicz P1,...,P7 wedlug podanych wzorow
        kazdy iloczyn przez Strassen
    oblicz C11,...,C22 i polacz bloki
    M = suma licznikow mnozen siedmiu wywolan
    S = suma licznikow dodawan + 18 * (n/2)^2
    return C, M, S
\end{lstlisting}

\begin{lstlisting}[language=Python,caption={Fragment implementacji siedmiu iloczynów.}]
P1, mul1, add1 = strassen(A11 + A22, B11 + B22)
P2, mul2, add2 = strassen(A21 + A22, B11)
P3, mul3, add3 = strassen(A11, B12 - B22)
P4, mul4, add4 = strassen(A22, B21 - B11)
P5, mul5, add5 = strassen(A11 + A12, B22)
P6, mul6, add6 = strassen(A21 - A11, B11 + B12)
P7, mul7, add7 = strassen(A12 - A22, B21 + B22)
C11 = P1 + P4 - P5 + P7
C22 = P1 - P2 + P3 + P6
\end{lstlisting}

\subsection{Nieparzyste rozmiary bez dopełniania zerami}
Dla nieparzystego $n>1$ oddzielamy ostatni wiersz i kolumnę:
\[
A=\begin{bmatrix}A_0&a\\r&\alpha\end{bmatrix},\qquad
B=\begin{bmatrix}B_0&b\\s&\beta\end{bmatrix}.
\]
Macierze $A_0,B_0$ mają parzysty rozmiar $(n-1)\times(n-1)$.
Tylko iloczyn $A_0B_0$ liczymy rekurencyjnie metodą Strassena.
Iloczyny z udziałem wektorów i skalarów wykonujemy bezpośrednio:
\[
AB=\begin{bmatrix}
 A_0B_0+as & A_0b+a\beta\\
 rB_0+\alpha s & rb+\alpha\beta
\end{bmatrix}.
\]

\begin{lstlisting}[caption={Pseudokod obsługi nieparzystego rozmiaru.}]
StrassenOdd(A, B):
    oddziel ostatni wiersz i kolumne A oraz B
    P = Strassen(A0, B0)
    C11 = P + Direct(a, s)
    C12 = Direct(A0, b) + Direct(a, beta)
    C21 = Direct(r, B0) + Direct(alpha, s)
    C22 = Direct(r, b) + Direct(alpha, beta)
    zsumuj liczniki i dodawania blokow
    return polacz(C11, C12, C21, C22), M, S
\end{lstlisting}

\subsection{Złożoność i liczba operacji}
Dla $n=2^j$ zachodzi
\[
 T(n)=7T(n/2)+\Theta(n^2),\qquad
 T(n)=\Theta(n^{\log_2 7})\approx\Theta(n^{2{,}8074}).
\]
Przy zejściu do skalarów liczba operacji jest równa
\[
 M=7^j,\qquad S=6(7^j-4^j).
\]
Dla nieparzystego $n=m+1$ dodatkowe iloczyny i sumy dają
\[
 M(n)=M(m)+3m^2+3m+1,\qquad
 S(n)=S(m)+3m^2+m.
\]
Koszt wydzielenia brzegu jest kwadratowy, więc nie zmienia dominującej
złożoności. Redukcja liczby mnożeń nie oznacza, że łączna liczba FLOPs
będzie mniejsza dla każdego badanego rozmiaru.

\section{Algorytm oparty na AlphaTensor}
\subsection{Faktoryzacja i sposób obliczeń}
Wykorzystujemy opublikowaną faktoryzację nad liczbami rzeczywistymi
dla iloczynu $4\times5$ przez $5\times5$, o randze 76\cite{alphatensor,data}.
Model AlphaTensor nie jest uruchamiany podczas mnożenia macierzy.
Program korzysta z zapisanych współczynników $U,V,W\in\{-1,0,1\}$.

Po podziale $A$ na 20 bloków, a $B$ na 25 bloków wyznaczamy
\[
 L_t=\sum_i U_{it}A_i,\quad R_t=\sum_j V_{jt}B_j,\quad
 H_t=L_tR_t\quad(t=1,\ldots,76),
\]
a następnie 20 bloków wyjściowych
$C_\ell=\sum_t W_{\ell t}H_t$.
Współczynniki zerowe pomijamy. Współczynnik $-1$ obsługujemy zmianą
znaku lub odejmowaniem, a nie dodatkowym mnożeniem przez $-1$.
Indeksy tablicy $W$ przestawiono tak, aby uzyskać wynik $AB$,
zamiast transponowanej konwencji danych źródłowych.

\begin{lstlisting}[caption={Pseudokod rekurencyjnego AlphaTensor.}]
AlphaTensor(A, B):
    (r, k) = rozmiar(A)
    s = liczba_kolumn(B)
    if r < 4 or k < 5 or s < 5:
        return Direct(A, B)
    dopelnij r do wielokrotnosci 4
    dopelnij k i s do wielokrotnosci 5
    podziel A na 4 x 5 blokow, B na 5 x 5 blokow
    for t = 1,...,76:
        L = kombinacja blokow A wedlug U[:,t]
        R = kombinacja blokow B wedlug V[:,t]
        H[t] = AlphaTensor(L, R)
    zloz 20 blokow wyniku wedlug W
    przytnij wynik do r x s
    return C, suma_mnozen, suma_dodawan_i_odejmowan
\end{lstlisting}

\begin{lstlisting}[language=Python,caption={Wykorzystanie współczynników w implementacji.}]
for left_terms, right_terms in zip(U, V):
    left, left_ops = combine(A_blocks, left_terms)
    right, right_ops = combine(B_blocks, right_terms)
    product, mul, add = alphatensor(left, right)
    products.append(product)
    multiplications += mul
    additions += left_ops + right_ops + add
\end{lstlisting}

\subsection{Dowolne rozmiary i interpretacja złożoności}
Dopełnianie jest wykonywane osobno na każdym poziomie rekurencji:
\[
r'=4\lceil r/4\rceil,\qquad k'=5\lceil k/5\rceil,
\qquad s'=5\lceil s/5\rceil.
\]
Dodane zera nie zmieniają właściwego iloczynu po przycięciu wyniku,
ale zwiększają koszt i powodują skoki liczby operacji.

Dla skalarnego jądra $4\times5$ przez $5\times5$ otrzymujemy
76 mnożeń i 549 binarnych dodawań lub odejmowań.
Zwykłe mnożenie tego samego przykładu wymaga 100 mnożeń i 80 dodawań.
Mniejsza liczba mnożeń nie oznacza zatem mniejszej sumy operacji.

Dla specjalnej rodziny wymiarów $4^j\times5^j$ przez $5^j\times5^j$
nie ma dopełniania i liczba mnożeń wynosi $76^j$.
Po podstawieniu $N=5^j$ daje to $N^{\log_5 76}\approx N^{2{,}691}$,
ale liczba wierszy wynosi wtedy $4^j$, a nie $N$.
Nie traktujemy tego wykładnika jako dowodu złożoności dla macierzy
kwadratowych $N\times N$. W ich przypadku głębokość zależy od
najszybciej malejących wymiarów dzielonych przez 5, a koszt liści
także od wymiaru wierszy dzielonego przez 4.
Wyniki kwadratowe oceniamy na podstawie rzeczywiście zliczonych operacji.

\section{Metodyka badań}
\subsection{Czas, pamięć i zakres rozmiarów}
Czas mierzymy funkcją \texttt{perf\_counter} wyłącznie wokół wywołania
algorytmu. Losowanie, wczytywanie, zapis oraz porównanie z Octave odbywają
się poza mierzonym odcinkiem. Dla kilku powtórzeń raportujemy medianę.
Nie wykonujemy pomiaru czasu równocześnie ze śledzeniem pamięci.
W AlphaTensor dane współczynników ładowane są przed pomiarem.

Historyczny zbiór Binéta zawiera rozmiary 1--128 oraz
160, 200, 256, 320, 384 i 512. Dla Strassena i AlphaTensor zbadano
rozmiary 1, 2, 3, 4, 5, 8, 17, 32, 64, 96, 128, 256 i 512.
Dla obu tych metod wykonywano 5 powtórzeń do 32, 3 do 128 i 1 powyżej 128.
W Binécie do 100 wykonywano 7 powtórzeń, do 200 po 3,
a dla większych rozmiarów po jednym.
Maksymalny zmierzony rozmiar wynosi 512; nie przedstawiamy
tych danych jako pełnego pomiaru wszystkich rozmiarów 1--1000.

Binét był pierwotnie mierzony na innych losowych parach i w osobnej sesji.
Strassen i AlphaTensor korzystały z tych samych zapisanych wejść.
Dlatego historyczne czasy pozwalają na orientacyjne porównanie,
a nie na kontrolowany pomiar przyspieszenia względem Binéta.
Nowe skrypty dla modułów Cython wczytują wspólne zapisane macierze
i zapisują wyniki w osobnych plikach z przyrostkiem \texttt{\_compiled}.

Pamięć badamy osobno dla rozmiarów
1, 2, 4, 8, 16, 32, 64, 128, 200 i 256.
Pole \texttt{peak\_traced\_memory\_mib} oznacza szczyt alokacji widocznych
dla \texttt{tracemalloc} podczas mnożenia. Nie jest to całkowita pamięć
procesu ani bezpośredni pomiar wszystkich alokacji natywnego kodu.
Macierze wejściowe i pamięć cache współczynników są tworzone wcześniej.
Pole \texttt{input\_buffers\_mib} opisuje osobno same bufory $A$ i $B$:
\[
 I=\frac{A.\texttt{nbytes}+B.\texttt{nbytes}}{2^{20}}
   =\frac{16n^2}{2^{20}}\ \text{MiB}
\]
dla dwóch macierzy kwadratowych typu float64.

\subsection{Weryfikacja}
Wyniki porównujemy element po elemencie z iloczynem \texttt{A*B}
w GNU Octave. Stosujemy warunek
\[
 |C_{ij}-C^{\mathrm{ref}}_{ij}|
 \le 10^{-12}+10^{-12}|C^{\mathrm{ref}}_{ij}|.
\]
Raportujemy także $E_{\max}=\max_{i,j}|C_{ij}-C^{\mathrm{ref}}_{ij}|$.
Bezwzględny błąd większy od $10^{-12}$ nie musi oznaczać niepowodzenia,
ponieważ warunek uwzględnia również tolerancję względną.

Wersje skompilowane porównano z oryginałami dla rozmiarów 1--33 i 64:
wyniki oraz liczniki były zgodne. Testy AlphaTensor obejmują dodatkowo
wszystkie 500 par macierzy bazowych jądra, rekurencyjny przykład
$16\times25$ przez $25\times25$ oraz przypadki prostokątne i wektorowe.

\section{Wyniki dotychczasowych pomiarów}
\subsection{Wspólne wykresy}
Poniższe wykresy przedstawiają \textbf{wersje Python przed kompilacją}.
Na wspólnych wykresach pokazano tylko rozmiary zmierzone dla wszystkich
trzech metod. Obie osie są logarytmiczne. Odcinki łączą próbki;
nie oznaczają pomiaru każdego rozmiaru pomiędzy nimi.
Przy zależności $t\approx cn^p$ wykres log--log ma postać
$\log t\approx\log c+p\log n$, czyli w przybliżeniu liniową.

\wykres{comparison_time.png}{Historyczne porównanie czasu wykonania trzech metod.}{fig:time}
\wykres{comparison_flops.png}{Liczba mnożeń oraz binarnych dodawań i odejmowań.}{fig:flops}
\wykres{comparison_memory.png}{Szczyt śledzonych alokacji; bufory wejściowe nie są uwzględnione.}{fig:memory}
\FloatBarrier

\subsection{Wybrane wartości}
\begin{table}[htbp]
\centering\small
\begin{tabular}{@{}rrrr@{}}
\toprule
$n$ & Binét [s] & Strassen [s] & AlphaTensor [s]\\
\midrule
128 & 7,080 & 4,023 & 6,712\\
256 & 56,590 & 29,380 & 12,820\\
512 & 429,224 & 210,828 & 357,719\\
\bottomrule
\end{tabular}
\caption{Historyczne czasy wersji Python. Dla 256 i 512 wykonano po jednym pomiarze.}
\end{table}

\begin{table}[htbp]
\centering\small
\begin{tabular}{@{}rrrr@{}}
\toprule
$n$ & Binét [FLOPs] & Strassen [FLOPs] & AlphaTensor [FLOPs]\\
\midrule
128 & 4\,177\,920 & 5\,666\,497 & 20\,204\,136\\
256 & 33\,488\,896 & 39\,960\,391 & 69\,267\,736\\
512 & 268\,173\,312 & 280\,902\,385 & 604\,962\,586\\
\bottomrule
\end{tabular}
\caption{Zliczone operacje arytmetyczne. Kompilacja nie zmienia tych liczników.}
\end{table}

\begin{table}[htbp]
\centering\small
\begin{tabular}{@{}lrr@{}}
\toprule
Metoda & Szczyt śledzonych alokacji [MiB] & Bufory wejściowe [MiB]\\
\midrule
Binét & 2,504 & 1,000\\
Strassen & 2,379 & 1,000\\
AlphaTensor & 3,749 & 1,000\\
\bottomrule
\end{tabular}
\caption{Historyczne pomiary pamięci dla $n=256$.}
\end{table}

W AlphaTensor liczba poziomów rekurencji i wielkość dopełnienia zależą
od wymiaru macierzy. To wyjaśnia nierówny wzrost kosztu i zmianę relacji
czasów pomiędzy 128, 256 i 512. Wyniku dla pojedynczego rozmiaru
nie można przenosić na wszystkie pozostałe.

Strassen dla $n=512$ wykonuje około 40,35 mln mnożeń, podczas gdy Binét
około 134,22 mln. Jednak dodatkowe sumy powodują, że Strassen ma tutaj
większą łączną liczbę operacji. Czasy zależą także od kosztu wywołań,
rozmiarów bloków i operacji NumPy, więc nie wynikają wyłącznie z sumy FLOPs.

\section{Przykład i porównanie z GNU Octave}
Dla zapisanego przypadku $4\times4$ otrzymano poniższe macierze.
Wartości w sprawozdaniu zaokrąglono do sześciu miejsc; porównanie
wykonano na pełnej precyzji plików CSV.
\[
A=\begin{bmatrix}
0.773956&0.438878&0.858598&0.697368\\
0.094177&0.975622&0.761140&0.786064\\
0.128114&0.450386&0.370798&0.926765\\
0.643865&0.822762&0.443414&0.227239
\end{bmatrix},
\]
\[
B=\begin{bmatrix}
0.554585&0.063817&0.827631&0.631664\\
0.758088&0.354526&0.970698&0.893121\\
0.778383&0.194639&0.466721&0.043804\\
0.154290&0.683049&0.744762&0.967510
\end{bmatrix}.
\]
Wynik referencyjny wynosi
\[
C^{\mathrm{ref}}=\begin{bmatrix}
1.537848&0.848438&1.986668&1.593172\\
1.505577&1.036961&1.965650&1.724703\\
0.844095&0.873047&1.406498&1.396071\\
1.361010&0.574301&1.707725&1.380811
\end{bmatrix}.
\]

\begin{lstlisting}[caption={Odtworzenie przykładu w GNU Octave.}]
A = dlmread('comparison/size_4/A.csv', ',');
B = dlmread('comparison/size_4/B.csv', ',');
C = A * B;
R = dlmread('comparison/size_4/C_strassen.csv', ',');
max_error = max(abs(C(:) - R(:)));
matches = all(abs(C(:) - R(:)) <= 1e-12 + 1e-12 * abs(C(:)));
\end{lstlisting}

\begin{table}[htbp]
\centering
\begin{tabular}{@{}lrr@{}}
\toprule
Metoda & $E_{\max}$, $n=4$ & $E_{\max}$, $n=512$\\
\midrule
Binét & $2{,}22\cdot10^{-16}$ & $1{,}14\cdot10^{-13}$\\
Strassen & $6{,}66\cdot10^{-16}$ & $4{,}39\cdot10^{-11}$\\
AlphaTensor & $2{,}22\cdot10^{-16}$ & $3{,}50\cdot10^{-12}$\\
\bottomrule
\end{tabular}
\caption{Maksymalny błąd bezwzględny wcześniejszych wyników względem Octave.
Wszystkie przypadki spełniają przyjęte kryterium tolerancji.}
\end{table}

Porównanie wykonano dla 13 rozmiarów od 1 do 512.
Dla $n=4$ AlphaTensor używa jeszcze mnożenia bezpośredniego;
jądro 76 iloczynów uruchamia się dopiero po spełnieniu progów wymiarów.
Weryfikacja poprawności obejmuje zatem także większe macierze.

\section{Wnioski i dalsze pomiary}
Binét daje prostą implementację o dokładnie znanej liczbie operacji
i złożoności sześciennej. Nierówne bloki pozwalają obsłużyć także
nieparzyste i prostokątne wymiary bez dopełniania.

Strassen zmniejsza liczbę mnożeń kosztem większej liczby sum i różnic.
Wydzielenie ostatniego wiersza i kolumny zapewnia obsługę nieparzystych
rozmiarów bez zmiany dominującej złożoności $\Theta(n^{\log_2 7})$.

AlphaTensor wykorzystuje 76 iloczynów zamiast 100 w jądrze prostokątnym.
W naszej implementacji koszt kombinacji liniowych i dopełniania
wpływa istotnie na czas, liczbę FLOPs i pamięć.
Metoda nie jest najszybsza dla każdego zmierzonego rozmiaru.

Przed końcowym porównaniem wersji skompilowanych należy wykonać wspólny
przebieg na zapisanych macierzach, uzupełnić zakres pomiarów i zastąpić
historyczne wykresy. Dopiero te dane pozwolą ocenić wpływ kompilacji
bez mieszania obu wariantów programu.

\begin{thebibliography}{9}
\bibitem{alphatensor}
A. Fawzi i in., \emph{Discovering faster matrix multiplication algorithms
with reinforcement learning}, Nature 610, 47--53, 2022.
\url{https://doi.org/10.1038/s41586-022-05172-4}.
\bibitem{data}
Google DeepMind, \emph{AlphaTensor: factorizations over the reals}.
\url{https://github.com/google-deepmind/alphatensor}.
Wykorzystano klucz $(4,5,5)$ z pliku \texttt{factorizations\_r.npz},
commit \texttt{1949163da3bef7e3eb268a3ac015fd1c2dbfc767}.
Dane: CC BY 4.0; wyodrębniono tablice i zmieniono indeksowanie $W$.
\bibitem{cython}
Dokumentacja Cython, \emph{Building Cython code}.
\url{https://docs.cython.org/en/latest/src/quickstart/build.html}.
\bibitem{repo}
Kod i dane projektu.
\url{https://github.com/lazorenkoo/Matrix-algorithms}.
\end{thebibliography}

\end{document}
